
\documentclass{article}
\usepackage{colortbl}
\usepackage{makecell}
\usepackage{multirow}
\usepackage{supertabular}

\begin{document}

\newcounter{utterance}

\centering \large Interaction Transcript for game `wordle', experiment `medium\_frequency\_words\_no\_clue\_no\_critic', episode 10 with nemotron{-}3{-}ultra.
\vspace{24pt}

{ \footnotesize  \setcounter{utterance}{1}
\setlength{\tabcolsep}{0pt}
\begin{supertabular}{c@{$\;$}|p{.15\linewidth}@{}p{.15\linewidth}p{.15\linewidth}p{.15\linewidth}p{.15\linewidth}p{.15\linewidth}}
    \# & \multicolumn{2}{c}{Player} && \multicolumn{2}{c}{Game Master} \\
    \hline

    \theutterance \stepcounter{utterance}  
    & & & \multicolumn{4}{p{0.6\linewidth}}{
        \cellcolor[rgb]{0.9,0.9,0.9}{
            \makecell[{{p{\linewidth}}}]{
                \texttt{\tiny{[P1$\langle$GM]}}
                \texttt{Eres un mago del lenguaje al que le gusta adivinar palabras usando las reglas dadas.} \\
\\ 
\texttt{¡Bienvenido a Wordle! Tienes seis intentos para adivinar la palabra objetivo, una palabra válida en español de cinco letras minúsculas. Por favor, usa las etiquetas "explicación:" e "intento:" para proporcionar una explicación concisa para cada intento.} \\
\\ 
\texttt{Por ejemplo, si tu intento es "parte", tu respuesta debería ser} \\
\texttt{explicación: esta es una palabra común en español de cinco letras, y estoy comenzando mi intento con esta palabra.} \\
\texttt{intento: parte} \\
\\ 
\texttt{Después de cada intento, tu respuesta será validada y recibirás retroalimentación indicando qué letras son correctas (verde), qué letras son correctas pero están en la posición incorrecta (amarillo) y qué letras son incorrectas (rojo). Esta retroalimentación puede ser útil para determinar qué letras incluir o excluir en tu siguiente intento.} \\
\\ 
\texttt{Por ejemplo, la retroalimentación para "parte" podría ser:} \\
\texttt{retroalimentación\_intento: p<amarillo> a<amarillo> r<verde> t<amarillo> e<rojo>} \\
\\ 
\texttt{La explicación debe contener detalles sobre cómo se usa el retroalimentación\_intento para llegar a un nuevo intento.} \\
\\ 
\texttt{No generes ningún otro texto.} \\
\\ 
\texttt{Comencemos con tu primer intento.} \\
            }
        }
    }
     \\ \\

    \theutterance \stepcounter{utterance}  
    & \multicolumn{4}{p{0.6\linewidth}}{
        \cellcolor[rgb]{0.561,0.737,0.561}{
            \makecell[{{p{\linewidth}}}]{
                \texttt{\tiny{[P1$\rangle$GM]}}
                \texttt{explicación: esta es una palabra común en español de cinco letras con vocales frecuentes (a, e) y consonantes habituales (r, s, t), ideal para iniciar la partida y obtener información máxima.} \\
\texttt{intento: rates} \\
            }
        }
    }
    & & \\ \\

    \theutterance \stepcounter{utterance}  
    & & & \multicolumn{2}{p{0.3\linewidth}}{
        \cellcolor[rgb]{0.9,0.9,0.9}{
            \makecell[{{p{\linewidth}}}]{
                \texttt{\tiny{[GM$|$GM]}}
                \texttt{Error: The guessed word is not a valid word for this game.} \\
            }
        }
    }
    & & \\ \\

    \theutterance \stepcounter{utterance}  
    & & & \multicolumn{4}{p{0.6\linewidth}}{
        \cellcolor[rgb]{0.9,0.9,0.9}{
            \makecell[{{p{\linewidth}}}]{
                \texttt{\tiny{[P1$\langle$GM]}}
                \texttt{Tu intento no es una palabra válida para este juego. Por favor, inténtalo de nuevo.} \\
\\ 
\texttt{Proporciona tu respuesta solo en este formato.} \\
\texttt{explicación: detalles} \\
\texttt{intento: palabra} \\
            }
        }
    }
     \\ \\

    \theutterance \stepcounter{utterance}  
    & \multicolumn{4}{p{0.6\linewidth}}{
        \cellcolor[rgb]{0.561,0.737,0.561}{
            \makecell[{{p{\linewidth}}}]{
                \texttt{\tiny{[P1$\rangle$GM]}}
                \texttt{explicación: "rates" no es válida; uso "ratés" sin tilde no existe, así que pruebo "ratos", palabra común con las mismas consonantes y vocales frecuentes para abrir el juego.} \\
\texttt{intento: ratos} \\
            }
        }
    }
    & & \\ \\

    \theutterance \stepcounter{utterance}  
    & & & \multicolumn{2}{p{0.3\linewidth}}{
        \cellcolor[rgb]{0.9,0.9,0.9}{
            \makecell[{{p{\linewidth}}}]{
                \texttt{\tiny{[GM$|$GM]}}
                \texttt{attempts = 1} \\
\texttt{target\_word = virus} \\
\texttt{guess = ratos} \\
\texttt{guess\_feedback = r<yellow> a<red> t<red> o<red> s<green>} \\
            }
        }
    }
    & & \\ \\

    \theutterance \stepcounter{utterance}  
    & & & \multicolumn{4}{p{0.6\linewidth}}{
        \cellcolor[rgb]{0.9,0.9,0.9}{
            \makecell[{{p{\linewidth}}}]{
                \texttt{\tiny{[P1$\langle$GM]}}
                \texttt{retroalimentación\_intento: r<amarillo> a<rojo> t<rojo> o<rojo> s<verde>} \\
\\ 
\texttt{Proporciona tu respuesta solo en este formato.} \\
\texttt{explicación: detalles} \\
\texttt{intento: palabra} \\
            }
        }
    }
     \\ \\

    \theutterance \stepcounter{utterance}  
    & \multicolumn{4}{p{0.6\linewidth}}{
        \cellcolor[rgb]{0.561,0.737,0.561}{
            \makecell[{{p{\linewidth}}}]{
                \texttt{\tiny{[P1$\rangle$GM]}}
                \texttt{explicación: La 's' está confirmada en la posición 5 (verde). La 'r' está en la palabra pero no en la posición 1 (amarillo), por lo que debe estar en la 2, 3 o 4. Las letras 'a', 't', 'o' se descartan (rojo). Probaré una palabra que termine en 's', contenga 'r' en posición distinta a la 1, y use nuevas vocales (e, i, u) y consonantes (l, n).} \\
\texttt{intento: rines} \\
            }
        }
    }
    & & \\ \\

    \theutterance \stepcounter{utterance}  
    & & & \multicolumn{2}{p{0.3\linewidth}}{
        \cellcolor[rgb]{0.9,0.9,0.9}{
            \makecell[{{p{\linewidth}}}]{
                \texttt{\tiny{[GM$|$GM]}}
                \texttt{Error: The guessed word is not a valid word for this game.} \\
            }
        }
    }
    & & \\ \\

    \theutterance \stepcounter{utterance}  
    & & & \multicolumn{4}{p{0.6\linewidth}}{
        \cellcolor[rgb]{0.9,0.9,0.9}{
            \makecell[{{p{\linewidth}}}]{
                \texttt{\tiny{[P1$\langle$GM]}}
                \texttt{Tu intento no es una palabra válida para este juego. Por favor, inténtalo de nuevo.} \\
\\ 
\texttt{Proporciona tu respuesta solo en este formato.} \\
\texttt{explicación: detalles} \\
\texttt{intento: palabra} \\
            }
        }
    }
     \\ \\

    \theutterance \stepcounter{utterance}  
    & \multicolumn{4}{p{0.6\linewidth}}{
        \cellcolor[rgb]{0.561,0.737,0.561}{
            \makecell[{{p{\linewidth}}}]{
                \texttt{\tiny{[P1$\rangle$GM]}}
                \texttt{explicación: "rines" no es válida en el diccionario del juego. Mantengo la 's' final (verde) y la 'r' en posición 2 (amarillo en intento anterior, probando posición distinta a la 1). Uso vocales nuevas 'e', 'i' y consonantes 'n', 'd' para formar "riend", pero no existe. Pruebo "riesg" inválido. Cambio estrategia: coloco 'r' en posición 3 con "seres", válida, 's' al inicio y final, 'e' vocal nueva.} \\
\texttt{intento: seres} \\
            }
        }
    }
    & & \\ \\

    \theutterance \stepcounter{utterance}  
    & & & \multicolumn{2}{p{0.3\linewidth}}{
        \cellcolor[rgb]{0.9,0.9,0.9}{
            \makecell[{{p{\linewidth}}}]{
                \texttt{\tiny{[GM$|$GM]}}
                \texttt{attempts = 2} \\
\texttt{target\_word = virus} \\
\texttt{guess = seres} \\
\texttt{guess\_feedback = s<red> e<red> r<green> e<red> s<green>} \\
            }
        }
    }
    & & \\ \\

    \theutterance \stepcounter{utterance}  
    & & & \multicolumn{4}{p{0.6\linewidth}}{
        \cellcolor[rgb]{0.9,0.9,0.9}{
            \makecell[{{p{\linewidth}}}]{
                \texttt{\tiny{[P1$\langle$GM]}}
                \texttt{retroalimentación\_intento: s<rojo> e<rojo> r<verde> e<rojo> s<verde>} \\
\\ 
\texttt{Proporciona tu respuesta solo en este formato.} \\
\texttt{explicación: detalles} \\
\texttt{intento: palabra} \\
            }
        }
    }
     \\ \\

    \theutterance \stepcounter{utterance}  
    & \multicolumn{4}{p{0.6\linewidth}}{
        \cellcolor[rgb]{0.561,0.737,0.561}{
            \makecell[{{p{\linewidth}}}]{
                \texttt{\tiny{[P1$\rangle$GM]}}
                \texttt{explicación: La 'r' está confirmada en la posición 3 (verde) y la 's' en la 5 (verde). Las letras 's' (pos 1) y 'e' se descartan (rojo). La 'r' no puede estar en posición 1 (amarillo en primer intento). Estructura: \_ \_ r \_ s. Letras descartadas: a, t, o, s, e. Vocales restantes: i, u. Consonantes probables: l, n, d, c, m, p, b, g, v. Pruebo "turns" (inválido), "tirns" (inválido). "Miras" usa 'a' (rojo). "Virar" no termina en 's'. "Durás" (verbo) o "dunas" (sustantivo). "Dunas" estructura: d{-}u{-}n{-}a{-}s (a rojo). "Dules" (e rojo). "Furor" no. "Juras" (a rojo). "Kilos" (k raro). "Lunas" (a rojo). "Muros" (o rojo). "Nudos" (o rojo). "Puros" (o rojo). "Rulos" (r pos 1 no). "Sulos" (s pos 1 no). "Tutás" (t rojo). "Uncir" no. "Virus" (v{-}i{-}r{-}u{-}s): vocales i, u nuevas; consonantes v nuevas; estructura coincide.} \\
\texttt{intento: virus} \\
            }
        }
    }
    & & \\ \\

    \theutterance \stepcounter{utterance}  
    & & & \multicolumn{2}{p{0.3\linewidth}}{
        \cellcolor[rgb]{0.9,0.9,0.9}{
            \makecell[{{p{\linewidth}}}]{
                \texttt{\tiny{[GM$|$GM]}}
                \texttt{attempts = 3} \\
\texttt{target\_word = virus} \\
\texttt{guess = virus} \\
\texttt{guess\_feedback = v<green> i<green> r<green> u<green> s<green>} \\
            }
        }
    }
    & & \\ \\

    \theutterance \stepcounter{utterance}  
    & & & \multicolumn{2}{p{0.3\linewidth}}{
        \cellcolor[rgb]{0.9,0.9,0.9}{
            \makecell[{{p{\linewidth}}}]{
                \texttt{\tiny{[GM$|$GM]}}
                \texttt{game\_result = WIN} \\
            }
        }
    }
    & & \\ \\

\end{supertabular}
}

\end{document}
